
\begin{frame}[fragile]{SQL, quantificateurs et négation}

Cette session présente les quantifications 
\alert{existentielle} et \alert{universelle}.

\vfill

Les quantificateurs permettent l'expression de la négation: ``je veux tous ces nuplets \alert{sauf} ceux-là''.

\vfill
Dans cette session:
\vfill

\begin{redbullet}
  \item Le quantificateur \texttt{exists} et l'expression du ``pour tout''
  \item La négation
  \item Equivalence des requêtes: plusieurs syntaxes, une seule signification.
\end{redbullet}

\vfill
\begin{blocgauche}
\alert{Ces diapositives correspondent au support en ligne disponible sur
le site \url{http://sql.bdpedia.fr/}}
\end{blocgauche}

\end{frame}

%%%%%%ù
\begin{frame}[fragile]{Le quantificateur ``exists``}

Reprenons la requête ``les logements où l’on peut faire du ski''.

\begin{minted}[baselinestretch=.9,
               fontsize=\small,
               framesep=2mm]{sql}
select distinct l.nom
from Logement as l, Activité as a
where l.code = a.codeLogement
and   a.codeActivité = 'Ski'
\end{minted}

\vfill

``a``  n'intervient pas dans le nuplet-résultat. 
On peut la remplacer par une variable liée.


\begin{minted}[baselinestretch=.9,
               fontsize=\small,
               framesep=2mm]{sql}
select distinct l.nom
from Logement as l
where exists (select ''
              from Activité as a
              where l.code = a.codeLogement
              and a.codeActivité = 'Ski')
\end{minted}


\vfill
\begin{blocgauche}
Légère reformulation: maintenant on cherche les logements
tels \alert{qu'il existe} une activité ``Ski''.
\end{blocgauche}
\end{frame}

%%%%%%ù
\begin{frame}[fragile]{Construction de formules complexes}

On peut construire des formules imbriquées (sous-requêtes SQL) 
sans limitation de profondeur.

\vfill 

Qui est allé dans les Alpes?
\begin{minted}[baselinestretch=.9,
               fontsize=\small,
               framesep=2mm]{sql}
     select distinct v.prénom, v.nom
     from Voyageur as v, Séjour as s, Logement as l
     where v.idVoyageur=s.idVoyageur
     and   s.codeLogement = l.code
     and   l.lieu = 'Alpes'
\end{minted}


\vfill
\begin{blocgauche}
Ni \texttt{s} ni \texttt{l} ne sont utilisées dans la construction du nuplet-résultat 
\end{blocgauche}
\end{frame}



%%%%%%ù
\begin{frame}[fragile]{Qui est allé dans les Alpes}


Avec quantificateur existentiel: "Les voyageurs tels \alert{qu'il existe} un de leurs séjours dans les Alpes".

\vfill 

\begin{minted}[baselinestretch=.9,
               fontsize=\small,
               framesep=2mm]{sql}
select distinct v.prénom, v.nom
from Voyageur as v
where exists (select '' 
              from Séjour as s, Logement as l
              where v.idVoyageur=s.idVoyageur
              and   s.codeLogement = l.code
              and   l.lieu = 'Alpes')
\end{minted}


\vfill
\begin{blocgauche}
Plus clair? Moins clair?
\end{blocgauche}
\end{frame}

%%%%%%ù
\begin{frame}[fragile]{Imbrication d'imbrication}

Avec quantificateur existentiel: "Les voyageurs tels qu'il \alert{existe} un de leurs séjours 
tel que son logement \alert{existe} dans les Alpes".

\vfill 

\begin{minted}[baselinestretch=.9,
               fontsize=\small,
               framesep=2mm]{sql}
select distinct v.prénom, v.nom
from Voyageur as v
where exists (select '' 
              from Séjour as s
              where v. idVoyageur=s.idVoyageur
              and exists (select '' 
                          from  Logement as l
                          where s.codeLogement = l.code
                          and   l.lieu = 'Alpes')
                    )
\end{minted}

\vfill
\begin{blocgauche}
Pas très naturel.
\end{blocgauche}
\end{frame}

%%%%%%ù
\begin{frame}[fragile]{Quantificateurs et négation}

Les logements qui \alert{ne proposent pas} de Ski.

\vfill 

\begin{minted}[baselinestretch=.9,
               fontsize=\small,
               framesep=2mm]{sql}
select distinct l.nom
from Logement as l
where not exists (select ''
                 from Activité as a
                 where l.code = a.codeLogement
                 and a.codeActivité = 'Ski')
\end{minted}

\vfill
\begin{blocgauche}
Correspond à la formulation: ``Les logements tels \alert{qu'il n'existe pas} d'activité Ski''.
\end{blocgauche}
\end{frame}



%%%%%%ù
\begin{frame}[fragile]{Quantificateur universel}

Les voyageurs qui sont allés dans \alert{tous} les logements

\vfill 
\begin{minted}[baselinestretch=.9,
               fontsize=\small,
               framesep=2mm]{sql}
     select distinct v.prénom, v.nom
     from Voyageur as v
     where  not exists (select ''
                        from Logement as l
                        where not exists  (select ''
                                           from Séjour as s
                                           where l.code = s.codeLogement
                                           and   v.idVoyageur = s.idVoyageur)
                        )
\end{minted}
\vfill
\begin{blocgauche}
Reformulation avec double négation: on cherche les voyageurs tels \alert{qu’il n’existe pas} 
de logement où \alert{ils ne sont pas} allés.
\end{blocgauche}
\end{frame}
%%%%%%%%%
\begin{frame}{À retenir}

SQL = un langage à la définition très précise.
\vfill

\begin{redbullet}
  \item  \alert{Tout ce qui peut s'exprimer par une formule logique est exprimable en SQL}. Ni plus, ni moins.

  \item \alert{Inversement}, tout ce qui ne s'exprime pas par une formule (boucles, incrémentations, etc.)
      ne s'exprime pas en SQL.
\end{redbullet}

\vfill
\begin{blocgauche}
Maîtriser SQL = savoir exprimer sa requête de manière rigoureuse. 
\end{blocgauche}

\vspace*{1cm}

Vous savez maintenant exprimer \alert{tout} ce qui est exprimable
en SQL (sauf agrégations)
\end{frame}



